%% =====================================================================
%% KSA / KFU - CCSIT Graduation Project Report Template (LaTeX2e)
%% THIS IS THE FILE YOU EDIT.
%% ---------------------------------------------------------------------
%% Files in this package:
%%   ccsit-preamble.tex   fixed layout            -- do not edit
%%   ccsit-macros.tex     your project details    -- EDIT THIS
%%   cover.tex            title page
%%   acknowledgment.tex   your acknowledgment     -- EDIT THIS
%%   undertaking.tex      signed declaration      -- do not edit the text
%%   bibliography.bib     your references         -- EDIT THIS
%%   CCSIT-Logo.png       college logo
%%   figures/             your images
%%
%% Compile with pdfLaTeX; Biber runs between the LaTeX passes:
%%     pdflatex main.tex / biber main / pdflatex main.tex / pdflatex main.tex
%% =====================================================================

\documentclass[11pt,a4paper]{report}

\input{ccsit-preamble}
%% =====================================================================
%%                 CCSIT / KFU MACROS (EDIT THESE)
%% =====================================================================
\newcommand{\ksa}{Kingdom of Saudi Arabia}
\newcommand{\moe}{Ministry of Education}
\newcommand{\kfu}{King Faisal University}
\newcommand{\ccsitCollege}{CCSIT, King Faisal University}

\newcommand{\ccsitFull}{College of Computer Sciences \& Information Technology}

\newcommand{\gpTitle}{AI Based Research Documents Repository Using Local Language Models}
\newcommand{\gpTitleShort}{Multimodal Document Repository} % used in header (keep short)
\newcommand{\gpDegree}{Bachelor of Science in Computer Science}
\newcommand{\gpDepartment}{Department of CS}

\newcommand{\gpSemester}{Second Semester}
\newcommand{\gpAcademicYear}{2025--2026}

%% Automatic Month Year based on compile date
\newcommand{\gpSubmissionMonth}{%
\ifcase\month\or
January\or February\or March\or April\or May\or June\or
July\or August\or September\or October\or November\or December%
\fi\space \number\year}

%% Students. For a team of fewer than four, do NOT delete or blank these definitions. Delete the matching lines from the "Prepared by" table in
%% cover.tex and from the signature table in main.tex instead.
\newcommand{\gpStudentA}{Raihanah Nezar Al Ali}
\newcommand{\gpStudentAID}{223007420}
\newcommand{\gpStudentB}{Maryam Ali Alkathim}
\newcommand{\gpStudentBID}{223021603}
\newcommand{\gpStudentC}{Fadak Hani Alibrahim}
\newcommand{\gpStudentCID}{223022618}

%% Supervision / Committee
\newcommand{\gpSupervisor}{Dr. Nora Abdullah Alkhaldi}
\newcommand{\gpCommitteeA}{Dr. Hala Hamdoun}
%% ----------------------------- Acknowledgment Text -----------------------------
\newcommand{\gpAcknowledgmentText}{
We would like to express our sincere gratitude to our supervisor, \textbf{\gpSupervisor}, for the continuous
guidance, valuable feedback, and support throughout the development of this graduation project. We
also extend our appreciation to the committee members and all faculty and staff in \ccsitFull\ for
their assistance and encouragement.
}
\input{undertaking}

%% =====================================================================
%%                           DOCUMENT
%% =====================================================================
\begin{document}

%% ----------------------------- Title Page -----------------------------
%% =====================================================================
%% cover.tex  --  the title page
%% ---------------------------------------------------------------------
%% Built entirely from the macros in ccsit-macros.tex. To change what
%% appears here, edit those macros.
%%
%% Two exceptions you may edit in this file:
%%   - teams of fewer than four students: delete the unused lines from
%%     the "Prepared by" table (and from undertaking.tex);
%%   - to show the co-supervisor or committee members 2 and 3, remove
%%     the leading % from those lines.
%% =====================================================================

\begin{titlepage}
\thispagestyle{empty}
\begin{center}

\vspace{4mm}
\includegraphics[width=0.5\textwidth]{CCSIT-Logo.png}

\vspace{6mm}

{\Large \textbf{\ksa}}\\[2mm]
{\Large \textbf{\moe}}\\[2mm]
{\Large \textbf{\kfu}}\\[2mm]
{\Large \textbf{\ccsitCollege}}\\[10mm]
{\large \textbf{\ccsitFull}}\\[10mm]
\vfill

{\LARGE \textbf{\gpTitle}\par}
\vspace{6mm}

{\large\itshape
\setstretch{1.05}
A project submitted\par
in partial fulfillment of the requirements for the degree of\par
\gpDegree\par
}

\vfill

\textbf{Prepared by}\\[2mm]
\begin{tabular}{rl}
\gpStudentA & (\gpStudentAID) \\
\gpStudentB & (\gpStudentBID) \\
\gpStudentC & (\gpStudentCID) \\
\end{tabular}

\vfill

\textbf{Supervised by}\\[2mm]
\gpSupervisor\\

\vspace{6mm}

\textbf{Committee Members}\\[2mm]
\gpCommitteeA\\

\vfill

{\gpSubmissionMonth}

\end{center}
\end{titlepage}

%% ----------------------------- Acknowledgment Page -----------------------------
\cleardoublepage
\thispagestyle{empty}
\begin{center}
{\Large \textbf{Acknowledgment}}
\end{center}

\vspace{6mm}
\gpAcknowledgmentText

%% ----------------------------- Undertaking -----------------------------
\cleardoublepage
\thispagestyle{empty}
\begin{center}
{\Large \textbf{Undertaking}}
\end{center}

\vspace{6mm}
\gpUndertakingText

\vfill
\noindent \textbf{Prepared by:}\\[3mm]
\begin{tabular}{@{}p{0.62\textwidth}p{0.34\textwidth}@{}}
\gpStudentA\ (\gpStudentAID) & Signature:\ \rule{45mm}{0.4pt}\\[4mm]
\gpStudentB\ (\gpStudentBID) & Signature:\ \rule{45mm}{0.4pt}\\[4mm]
\gpStudentC\ (\gpStudentCID) & Signature:\ \rule{45mm}{0.4pt}\\[4mm]
\end{tabular}

\vspace{6mm}
\noindent Date:\ \rule{55mm}{0.4pt}

%% ----------------------------- Front Matter -----------------------------
\cleardoublepage
\pagenumbering{roman}
\setcounter{page}{1}

\section*{Abstract}
\addcontentsline{toc}{section}{Abstract}

Scientific research documents convey important data both in their written form and through various graphical features, such as intricate figures, tables, and charts. Yet most conventional \ac{RAG} systems focus mainly on text, leading to a substantial loss of the key information contained in these visual elements. Moreover, existing AI-based document processing systems place a great deal of reliance on cloud-based \ac{API}, thereby creating serious risks in the areas of data privacy, protection of intellectual property, and the general reliability of the system. To overcome these important limitations, this project presents a thorough and privacy-preserving solution that uses AI to create a local multi-modal repository of research documents. The system in question is intended to function completely offline, thus keeping sensitive academic and proprietary data secure while simultaneously carrying out a comprehensive analysis of both the text and the visual parts of scientific literature.\\

In technical terms, the system functions by using deep layout-aware parsing and carries out separate extraction procedures for the textual and visual elements to maintain the original document structure. For the purpose of information retrieval, the architecture makes use of an advanced multi-modal method that combines hybrid text search-this involving both keyword matching and dense semantic analysis-with independent visual search operations, the two different streams being smoothly integrated by means of the \ac{RRF} algorithm. To produce responses the system uses a Vision‑Language Model that is hosted locally and fine‑tuned on scientific multi-modal datasets. This approach removes the need for any cloud services completely. The project contains three components. First, a functional web‑based repository makes document management easy. Second a powerful local AI processing pipeline guarantees offline operation. Third a thorough performance evaluation checks both retrieval and question‑answering accuracy. Finally, the project delivers an intelligent environment that allows researchers and educational institutions to query and understand complex scientific literature efficiently while keeping full control of their sensitive data.


\vspace{4mm}
\noindent\textbf{Keywords:} Research documents; Vision Language Model (VLM); Retrieval Augmented Generation (RAG); Local model.

\cleardoublepage
\tableofcontents

\cleardoublepage
\listoftables
\addcontentsline{toc}{section}{List of Tables}




\cleardoublepage
\listoffigures
\addcontentsline{toc}{section}{List of Figures}

%% ----------------------------- Acronyms -----------------------------
\cleardoublepage
\section*{List of Acronyms}
\addcontentsline{toc}{section}{List of Acronyms}

\begin{acronym}[123456]  
  \acro{LLMs}{Large Language Models}
  \acro{VLMs}{Vision Language Models}
  \acro{RAG}{Retrieval Augmented Generation} 
  \acro{IR}{Information Retrieval} 
  \acro{API}{Application Programming Interface}
  \acro{RRF}{Reciprocal Rank Fusion}  
\end{acronym}

%% ----------------------------- Main Matter -----------------------------
\cleardoublepage
\pagenumbering{arabic}
\setcounter{page}{1}

\section{Introduction}
Research documents contain valuable information in both textual and visual forms, and nowadays, the use of \ac{RAG} systems is increasing rapidly. However, conventional \ac{RAG} systems primarily focus on understanding and retrieving textual information \cite{Tanaka2025VDocRAGRG}, resulting in loss of significant knowledge displayed through visual content, such as figures, diagrams, tables, and charts, in the retrieval process. This limitation affects the completeness of retrieved information and the quality of the generated results.\\

To address this limitation, our project aims to apply local multi-modal retrieval techniques that ensure visual content is considered in information retrieval, which will support a wider range of research document questions and improve the quality of generated results. To provide a clear overview of the project, we will start by discussing fundamental components relevant to the project followed by its core objectives. Then, we will explain the importance of the project, which will provide the context for the specific problem we want to address. The last aspect to be discussed in this section is the project boundaries and expected outcomes.

\subsection{Background}
Processing multi-modal documents requires specialized pipeline to enhance the efficiency of the use of recent technologies in this field. The pipeline mainly relies on four interconnected techniques that are fundamentally relevant to context search and answer generation, which are \ac{IR}, \ac{LLMs}, \ac{VLMs}, and Retrieval Augmented Generation (\ac{RAG}). 

\subsubsection{Information Retrieval}
\subsubsection{Large Language Models}
\subsubsection{Vision Language Models}
\subsubsection{Retrieval Augmented Generation}
\subsection{objectives}
\subsection{Motivation}
\subsection{Problem Statement}
\subsection{Scope and Deliverables}
\subsubsection{Definition of the Project's Scope}
\subsubsection{Alignment with KFU Identity and Society Impact}
\subsubsection{Expected Project Deliverables}

\section{Literature Review}

\section{Methodology}

\subsection{Challenges and Proposed Strategies}
\subsubsection{Challenge 1: Limited Retrieval of Textless Visual Information}
\subsubsection{Challenge 2: High Computational Cost}
\subsubsection{Challenge 3: Limited Multi-Modal Information Extraction}
\subsubsection{Challenge 4: Limitations in Context Retrieval}
\subsubsection{Challenge 5: Limitations in Cross-Modal Retrieval}

\subsection{Justifications of the Selected Approach}
\subsubsection{Alternative Multi-Modal \ac{RAG} Pipeline}
\subsubsection{Proposed Multi-Modal \ac{RAG} Pipeline}

\subsection{Discussion of Tools and Techniques to be Used}

\section{Project Requirements}
\subsection{Functional Requirements}
\subsection{Non-Functional Requirements}
\subsubsection{Hardware Requirements}
\subsubsection{Software Requirements}

\section{Project Timeline}
Hello World Team! This project follows software engineering principles \cite{Gosmar2026MADPAM}.
Similar AI-based solutions have been reported in recent studies \cite{Dong2026DocResearcherAU}.
A comparable system design was presented in \cite{Smith2007AnOO}.

Acronyms are expanded automatically on first use. 
The proposed system exposes an API that applies
AI and NLP techniques on a GPU, and every later mention of the API prints the
short form only.

%% ----------------------------- References -----------------------------
\cleardoublepage
\printbibliography
\addcontentsline{toc}{section}{References}

%% ----------------------------- Appendices -----------------------------
% \appendix
% \cleardoublepage
%% ----------------------------- AppendixA -----------------------------
% \section*{Appendix A: Additional Materials}
% \addcontentsline{toc}{section}{Appendix A: Additional Materials}

% The following equation is provided only as an illustrative example to demonstrate how to include a general numbered equation in this report template.

% \begin{equation}
% \label{eq:general-example}
% y = ax + b
% \end{equation}

% \begin{subequations}
% \label{eq:subequations-example}
% \begin{align}
% y &= ax + b \label{eq:subeq-a}\\
% x &= \frac{y-b}{a} \label{eq:subeq-b}\\
% x^{10} &= y - b \label{eq:subeq-c}\\
% f^{(k)}(x) &\neq f_i(x)
% \end{align}
% \end{subequations}

% Table~\ref{tab:general-comparison} is included only as an example to demonstrate table formatting in this report template.

% \begin{table}[H]
% \centering
% \caption{General comparison between existing solutions and the proposed system.}
% \label{tab:general-comparison}
% \begin{tabular}{|l|c|c|}
% \hline
% \textbf{Criteria} & \textbf{Existing Solutions} & \textbf{Proposed System} \\
% \hline
% Ease of use & Moderate & High \\
% \hline
% Accuracy of results & Moderate & Improved \\
% \hline
% User interface clarity & Basic & Enhanced \\
% \hline
% Automation level & Limited & High \\
% \hline
% Reliability & Moderate & High \\
% \hline
% Scalability & Limited & Good \\
% \hline
% Reporting and documentation & Basic & Detailed \\
% \hline
% \end{tabular}
% \end{table}

% The following table is provided only as an illustrative example to demonstrate how to include a bordered table with \textbf{merged column headers} in this report template (Table~\ref{tab:general-comparison-merged}).

% \begin{table}[H]
% \centering
% \caption{Comparison of system characteristics.}
% \label{tab:general-comparison-merged}
% \begin{tabular}{|l|c|c|}
% \hline
% \textbf{Feature} & \multicolumn{2}{c|}{\textbf{System Type}} \\
% \cline{2-3}
%  & \textbf{Traditional Approach} & \textbf{Proposed Approach} \\
% \hline
% User experience & Average & Improved \\
% \hline
% System flexibility & Low & High \\
% \hline
% Maintenance effort & High & Lower \\
% \hline
% Processing speed & Moderate & Faster \\
% \hline
% Overall efficiency & Moderate & High \\
% \hline
% \end{tabular}
% \end{table}

% The following table is provided only as an illustrative example to demonstrate how to include a bordered table with \textbf{merged rows} in this report template (Table~\ref{tab:general-comparison-categories}).

% \begin{table}[H]
% \centering
% \caption{General comparison of project features.}
% \label{tab:general-comparison-categories}
% \begin{tabular}{|l|l|c|c|}
% \hline
% \textbf{Category} & \textbf{Criteria} & \textbf{Existing} & \textbf{Proposed} \\
% \hline
% \multirow{2}{*}{Usability}
% & Ease of use & Medium & High \\
% \cline{2-4}
% & Accessibility & Medium & High \\
% \hline
% \multirow{2}{*}{Performance}
% & Reliability & Medium & High \\
% \cline{2-4}
% & Response speed & Medium & High \\
% \hline
% \multirow{2}{*}{Output Quality}
% & Report clarity & Basic & Detailed \\
% \cline{2-4}
% & Output consistency & Medium & High \\
% \hline
% \end{tabular}
% \end{table}

% The following table is provided only as an illustrative example to demonstrate how to include a bordered table with \textbf{merged rows and merged column headers} in this report template (Table~\ref{tab:merged-rows-cols-example}).

% \begin{table}[H]
% \centering
% \caption{Illustrative comparison table (example only).}
% \label{tab:merged-rows-cols-example}
% \begin{tabular}{|l|l|c|c|}
% \hline
% \multirow{2}{*}{\textbf{Category}}
% & \multirow{2}{*}{\textbf{Criteria}}
% & \multicolumn{2}{c|}{\textbf{Comparison}} \\
% \cline{3-4}
%  &  & \textbf{Option A} & \textbf{Option B} \\
% \hline
% \multirow{2}{*}{Usability}
% & Ease of use & High & Medium \\
% \cline{2-4}
% & Clarity of interface & High & Medium \\
% \hline
% \multirow{2}{*}{Quality}
% & Reliability & High & Medium \\
% \cline{2-4}
% & Output clarity & High & Medium \\
% \hline
% \multirow{2}{*}{Efficiency}
% & Time required & Low & Medium \\
% \cline{2-4}
% & Effort required & Low & Medium \\
% \hline
% \end{tabular}
% \end{table}

% The following code listing is provided only as an illustrative example to demonstrate how to include a numbered and captioned code excerpt in this report template (Listing~\ref{lst:example}).

% \begin{lstlisting}[language=Python, caption={Example function that normalizes an input string.}, label={lst:example}]
% def normalize(text):
%     """Return the text in lower case, without leading or trailing spaces."""
%     return text.strip().lower()
% \end{lstlisting}

% As shown in Listing~\ref{lst:example}, source code is typeset with line numbers and a caption placed above the excerpt. Set \texttt{language} to the language of your own code, for example \texttt{Java}, \texttt{C}, \texttt{Python}, or \texttt{SQL}. Include only short, meaningful excerpts in the report; complete source files belong in the submitted code archive.

% %% ----------------------------- AppendixB -----------------------------
% \section*{Appendix B: Project Screenshots}
% \addcontentsline{toc}{section}{Appendix B: Project Screenshots}

% \begin{figure}[H]
% \centering
% \includegraphics[width=0.5\textwidth]{ex.png}
% \caption{Official CCSIT logo used in this graduation project report.}
% % \captionof{table}{Official CCSIT logo used in this graduation project report.}
% \label{fig:ccsit-logo-overall}
% \end{figure}

% As shown in Figure~\ref{fig:ccsit-logo-overall}, the official CCSIT logo is presented for institutional identification and report branding.

% \begin{figure}[H]
% \centering

% \begin{subfigure}[b]{0.45\textwidth}
% \centering
% \includegraphics[width=0.75\textwidth]{ex.png}
% \caption{King Faisal University (KFU) logo.}
% \label{fig:kfu-logo}
% \end{subfigure}
% \hfill
% \begin{subfigure}[b]{0.45\textwidth}
% \centering
% \includegraphics[width=0.75\textwidth]{ex.png}
% \caption{College of Computer Sciences \& Information Technology (CCSIT) logo.}
% \label{fig:ccsit-logo}
% \end{subfigure}

% \caption{Institutional logos used for official report branding.}
% \label{fig:logos}
% \end{figure}

% Figure~\ref{fig:logos} presents the institutional logos used in this report.

\end{document}